% part: intuitionistic-logic
% chapter: introduction
% section: natural-deduction

\documentclass[../../../include/open-logic-section]{subfiles}

\begin{document}

\olfileid{int}{int}{ntd}

\olsection{স্বাভাবিক অবরোহণ}

$\FalseCl$-এর নিয়মগুলি ছাড়া স্বাভাবিক অবরোহণ স্বজ্ঞাবাদী
যুক্তিবিদ্যার একটি প্রচলিত !!{derivation} ব্যবস্থা। নিয়মগুলি
এখানে আবার লিখে BHK ব্যাখ্যা দিয়ে তাদের প্রেরণা দেখাব। প্রতিটি
ক্ষেত্রে নিয়মটিকে এমনভাবে ভাবা যায়: পূর্বপক্ষগুলির নির্মাণ থাকলে
সিদ্ধান্তটিরও নির্মাণ পাওয়া যায়।

স্বাভাবিক অবরোহণের !!{derivation}s-এ অব্যাহতি না-পাওয়া অনুমান থাকে।
তাই~$\Gamma$ !!{undischarged} অনুমানগুলি থেকে $!A$-এর এমন
!!a{derivation}-কে এমন একটি অপেক্ষক হিসেবে ভাবা উচিত, যা
সব $!B \in \Gamma$-এর নির্মাণ থেকে~$!A$-এর নির্মাণ দেয়। কোনো
!!{undischarged} অনুমান ছাড়া~$!A$-এর !!a{derivation} থাকলে BHK
ব্যাখ্যার অর্থে~$!A$-এর একটি নির্মাণ আছে। তবে আলোচনার সুবিধায়
প্রয়োজন না হলে~$\Gamma$ উল্লেখ করব না।

শুধু অনুমান~$!A$-ই !!{undischarged} অনুমান~$!A$ থেকে~$!A$-এর
!!a{derivation}। এটি BHK ব্যাখ্যার সঙ্গে মেলে: নির্মাণগুলির উপর
অভেদ অপেক্ষক~$!A$-এর যেকোনো নির্মাণকে~$!A$-এর নির্মাণে পাঠায়।

\subsection{সংযোজন}

\begin{defish}
\AxiomC{$!A$}
\AxiomC{$!B$}
\RightLabel{\Intro{\land}}
\BinaryInfC{$!A \land !B$}
\DisplayProof
\hfill
\begin{tabular}{l}
\AxiomC{$!A \land !B$}
\RightLabel{\Elim{\land}}
\UnaryInfC{$!A$}
\DisplayProof
\\[3ex]
\AxiomC{$!A \land !B$}
\RightLabel{\Elim{\land}}
\UnaryInfC{$!B$}
\DisplayProof
\end{tabular}
\end{defish}

ধরো, যথাক্রমে $!A_1$ ও $!A_2$-এর নির্মাণ $N_1$, $N_2$ আছে। তাহলে
$!A_1 \land !A_2$-এরও একটি নির্মাণ আছে, অর্থাৎ জোড়াটি
$\tuple{N_1, N_2}$।

% BN-SRC-395: the pair eliminates A_1 and A_2, so the conjunction is A_1 and A_2.
BHK ব্যাখ্যায় $!A_1 \land !A_2$-এর একটি নির্মাণ হলো জোড়া
$\tuple{N_1, N_2}$। তাই এমন একটি জোড়া আছে ধরে নিলে প্রতিটি
সংযোজিত অংশের নির্মাণ পাই: $N_1$ হলো~$!A_1$-এর এবং $N_2$ হলো
$!A_2$-এর নির্মাণ।

\subsection{শর্তসূচক}

\begin{defish}
  \AxiomC{$\Discharge{!A}{u}$}
  \DeduceC{$!B$}
  \DischargeRule{$\Intro{\lif}$}{u}
  \UnaryInfC{$!A \lif !B$}
  \DisplayProof
\hfill
  \AxiomC{$!A \lif !B$}
  \AxiomC{$!A$}
  \RightLabel{$\Elim{\lif}$}
  \BinaryInfC{$!B$}
  \DisplayProof
\end{defish}

!!{undischarged} অনুমান~$!A$ থেকে~$!B$-এর !!a{derivation} থাকলে
এমন একটি অপেক্ষক~$f$ আছে, যা~$!A$-এর নির্মাণকে~$!B$-এর নির্মাণে
পাঠায়। সেই অপেক্ষকই $!A \lif !B$-এর একটি নির্মাণ। তাই $\Intro\lif$-এর
পূর্বপক্ষের নির্মাণ যদি~$!A$-এর নির্মাণের উপর শর্তাধীন হয়, তবে
সিদ্ধান্ত $!A \lif !B$-এরও নির্মাণ আছে।

অন্যদিকে, ধরো $!A$-এর নির্মাণ~$N$ এবং $!A \lif !B$-এর
নির্মাণ~$f$ আছে। $!A \lif !B$-এর নির্মাণ এমন একটি অপেক্ষক,
যা $!A$-এর নির্মাণকে~$!B$-এর নির্মাণে পাঠায়। ফলে $f(N)$ হলো
$!B$-এর একটি নির্মাণ; অর্থাৎ $\Elim\lif$-এর সিদ্ধান্তের নির্মাণ আছে।

\subsection{বিকল্প}

\begin{defish}
\begin{tabular}{l}
\AxiomC{$!A$}
\RightLabel{\Intro{\lor}}
\UnaryInfC{$!A \lor !B$}
\DisplayProof
\\[3ex]
\AxiomC{$!B$}
\RightLabel{\Intro{\lor}}
\UnaryInfC{$!A \lor !B$}
\DisplayProof
\end{tabular}
\hfill
\AxiomC{$!A \lor !B$}
\AxiomC{$\Discharge{!A}{n}$}
\DeduceC{$!C$}
\AxiomC{$\Discharge{!B}{n}$}
\DeduceC{$!C$}
\DischargeRule{\Elim{\lor}}{n}
\TrinaryInfC{$!C$}
\DisplayProof
\end{defish}

যদি~$!A_i$-এর একটি নির্মাণ~$N_i$ থাকে, তবে সেটি দিয়ে
$!A_1 \lor !A_2$-এর নির্মাণ~$\tuple{i, N_i}$ পাওয়া যায়।
অন্যদিকে, ধরো $!A_1 \lor !A_2$-এর একটি নির্মাণ আছে, অর্থাৎ একটি
জোড়া~$\tuple{i, N_i}$ যেখানে $N_i$ হলো~$!A_i$-এর নির্মাণ;
আর এমন অপেক্ষক $f_1$, $f_2$-ও আছে, যেগুলি যথাক্রমে $!A_1$, $!A_2$-এর
নির্মাণকে~$!C$-এর নির্মাণে পাঠায়। তাহলে $f_i(N_i)$ হলো~$!C$-এর
নির্মাণ, অর্থাৎ $\Elim\lor$-এর সিদ্ধান্তের নির্মাণ।

\subsection{অসংগতি}

\begin{defish}
  \AxiomC{$\lfalse$}
  \RightLabel{\FalseInt}
  \UnaryInfC{$!A$}
  \DisplayProof
\end{defish}

!!{undischarged} অনুমান~$!B_1$, \dots, $!B_n$ থেকে~$\lfalse$-এর
একটি !!a{derivation} থাকলে এমন একটি অপেক্ষক~$f(M_1,
\dots, M_n)$ আছে, যা $!B_1$, \dots, $!B_n$-এর নির্মাণকে
$\lfalse$-এর নির্মাণে পাঠায়। যেহেতু $\lfalse$-এর কোনো নির্মাণ
নেই, $!B_1$, \dots, $!B_n$ সকলের নির্মাণও একসঙ্গে থাকতে পারে না।
ফলে $f$-এর এই ধর্মও আছে: $M_1$, \dots, $M_n$ যথাক্রমে
$!B_1$, \dots, $!B_n$-এর নির্মাণ \emph{হলে}, $f(M_1, \dots, M_n)$
হবে~$!A$-এর একটি নির্মাণ।

\subsection{$\lnot$-এর নিয়ম}

যেহেতু $\lnot !A$-কে $!A \lif \lfalse$ হিসেবে সংজ্ঞায়িত করা
হয়েছে, কঠোর অর্থে~$\lnot$-এর জন্য আলাদা নিয়ম দরকার নেই। তবু
নিয়ম দিলে সেগুলি এমন হতো:

\begin{defish}
\AxiomC{$\Discharge{!A}{n}$}
\noLine
\DeduceC{$\lfalse$}
\DischargeRule{\Intro{\lnot}}{n}
\UnaryInfC{$\lnot !A$}
\DisplayProof
\hfill
\AxiomC{$\lnot !A$}
\AxiomC{$!A$}
\RightLabel{\Elim{\lnot}}
\BinaryInfC{$\lfalse$}
\DisplayProof
\end{defish}


\subsection{\usetoken{P}{derivation}-এর উদাহরণ}

\begin{enumerate}
\item $\Proves !A \lif (\lnot !A \lif \lfalse)$, অর্থাৎ $\Proves !A
  \lif ((!A \lif \lfalse) \lif \lfalse)$
  \begin{prooftree}
    \AxiomC{$\Discharge{!A}{2}$}
    \AxiomC{$\Discharge{!A \lif \lfalse}{1}$}
    \RightLabel{\Elim\lif}
    \BinaryInfC{$\lfalse$}
    \DischargeRule{\Intro\lif}{1}
    \UnaryInfC{$(!A \lif \lfalse)\lif \lfalse$}
    \DischargeRule{\Intro\lif}{2}
% BN-SRC-396: parenthesize the consequent to match the stated formula and tree.
    \UnaryInfC{$!A \lif ((!A \lif \lfalse) \lif \lfalse)$}
  \end{prooftree}

\item $\Proves ((!A \land !B) \lif !C) \lif (!A \lif (!B \lif !C))$
  \begin{prooftree}
    \AxiomC{$\Discharge{(!A \land !B) \lif !C}{3}$}
    \AxiomC{$\Discharge{!A}{2}$}
    \AxiomC{$\Discharge{!B}{1}$}
    \RightLabel{\Intro\land}
    \BinaryInfC{$!A \land !B$}
    \RightLabel{\Elim\lif}
    \BinaryInfC{$!C$}
    \DischargeRule{\Intro\lif}{1}
    \UnaryInfC{$!B \lif !C$}
    \DischargeRule{\Intro\lif}{2}
    \UnaryInfC{$!A \lif (!B \lif !C)$}
    \DischargeRule{\Intro\lif}{3}
    \UnaryInfC{$((!A \land !B) \lif !C) \lif (!A \lif (!B \lif !C))$}
  \end{prooftree}

\item $\Proves \lnot(!A \land \lnot !A)$, অর্থাৎ $\Proves (!A \land (!A
  \lif \lfalse))\lif \lfalse$
  \begin{prooftree}
    \AxiomC{$\Discharge{!A \land (!A \lif \lfalse)}{1}$}
    \RightLabel{\Elim\land}
    \UnaryInfC{$!A \lif \lfalse$}
    \AxiomC{$\Discharge{!A \land (!A \lif \lfalse)}{1}$}
    \RightLabel{\Elim\land}
    \UnaryInfC{$!A$}
    \RightLabel{\Elim\lif}
    \BinaryInfC{$\lfalse$}
    \DischargeRule{\Intro\lif}{1}
    \UnaryInfC{$(!A \land (!A \lif \lfalse))\lif \lfalse$}
  \end{prooftree}

\item $\Proves \lnot\lnot(!A \lor \lnot !A)$, অর্থাৎ $\Proves ((!A \lor
  (!A \lif \lfalse))\lif \lfalse)\lif \lfalse$
  \begin{prooftree}\hskip-3em
    \AxiomC{$\Discharge{(!A \lor (!A \lif \lfalse))\lif \lfalse}{2}$}
    \AxiomC{$\Discharge{(!A \lor (!A \lif \lfalse))\lif \lfalse}{2}$}
    \AxiomC{$\Discharge{!A}{1}$}
    \RightLabel{\Intro\lor}
    \UnaryInfC{$!A \lor (!A \lif \lfalse)$}
    \RightLabel{\Elim\lif}
    \BinaryInfC{$\lfalse$}
    \DischargeRule{\Intro\lif}{1}
    \UnaryInfC{$!A \lif \lfalse$}
    \RightLabel{\Intro\lor}
    \UnaryInfC{$!A \lor (!A \lif \lfalse)$}
    \RightLabel{\Elim\lif}
    \insertBetweenHyps{\hskip -4em}
    \BinaryInfC{$\lfalse$}
    \DischargeRule{\Intro\lif}{2}
    \UnaryInfC{$((!A \lor (!A \lif \lfalse))\lif \lfalse)\lif \lfalse$}
  \end{prooftree}
\end{enumerate}

\begin{prop}
  স্বজ্ঞাবাদী যুক্তিবিদ্যার স্বাভাবিক অবরোহণে $\Gamma \Proves !A$
  হলে ধ্রুপদি যুক্তিবিদ্যাতেও $\Gamma \Proves !A$। বিশেষত, $!A$
  স্বজ্ঞাবাদী উপপাদ্য হলে সেটি ধ্রুপদি উপপাদ্যও।
\end{prop}

\begin{proof}
  স্বাভাবিক অবরোহণের প্রতিটি নিয়ম ধ্রুপদি স্বাভাবিক অবরোহণেরও
  নিয়ম। তাই স্বজ্ঞাবাদী যুক্তিবিদ্যার প্রতিটি !!{derivation} ধ্রুপদি
  যুক্তিবিদ্যাতেও !!a{derivation}।
\end{proof}

\begin{prob}
  নিচের !!{formula}s-এর স্বজ্ঞাবাদী !!{derivation}s দাও:
  \begin{enumerate}
    \item $(\lnot !A \lor !B) \lif (!A \lif !B)$
    \item $\lnot\lnot\lnot !A \lif \lnot !A$
    \item $\lnot\lnot (!A \land !B) \liff (\lnot\lnot !A\land \lnot \lnot !B)$
    \item $\lnot (!A \lor !B) \liff (\lnot !A \land \lnot !B)$
    \item $(\lnot !A \lor \lnot !B) \lif \lnot (!A \land !B)$
    \item $\lnot\lnot ( !A \land !B) \lif  (\lnot\lnot !A \lor
    \lnot\lnot !B)$
  \end{enumerate}
\end{prob}

\end{document}
